\begin{abstract}
Although large unsupervised language models (LMs) acquire extensive knowledge about the world and display some reasoning capabilities, exercising fine-grained control over their outputs remains challenging due to the purely unsupervised nature of their pretraining. Common approaches to enhancing such controllability involve gathering human judgments regarding the relative merits of model outputs and subsequently fine-tuning the original LM to match these expressed preferences, frequently leveraging reinforcement learning from human feedback (RLHF). However, the RLHF pipeline is often intricate and prone to instability: it typically requires learning a separate reward model based on human preferences, followed by the application of reinforcement learning to adjust the pretrained LM to optimize this learned reward function, all while maintaining similarity to the base model. In this work, we propose a novel reparameterization of the reward model used in RLHF, which permits the closed-form derivation of the optimal policy. This leads to a reformulation of the standard RLHF objective as a straightforward classification loss. Our approach, termed \textit{Direct Preference Optimization} (DPO), is efficient, robust, and requires minimal computational resources, dispensing with the necessity for language model sampling during fine-tuning as well as extensive hyperparameter searches. Empirical results demonstrate that DPO can fine-tune LMs to reflect human preferences as effectively or better than existing techniques. Specifically, DPO outperforms PPO-based RLHF methods in steering generation sentiment, and achieves comparable or improved performance in summarization and single-turn conversational tasks, all with a much simpler implementation and training process.
\end{abstract}

\section{Introduction}
Large unsupervised language models (LMs), when exposed to vast amounts of data, develop a range of remarkable abilities~\citep{chowdhery2022palm, brown2020language, touvron2023llama,bubeck2023sparks}. Despite their powerful performance, these models are trained on human-generated data reflecting a broad spectrum of intentions, competencies, and objectives. Not all of these human-derived tendencies are desirable; for instance, an AI coding assistant should be able to \textit{recognize} prevalent programming errors in order to suggest corrections, yet during code generation, we aim for the model to emulate the (sometimes infrequent) top-tier coding practices within its training set. In a similar vein, we may expect a language model to be \textit{informed} about widely held misconceptions—such as those accepted by 50\% of individuals—while ensuring it refrains from propagating such misconceptions as factual in half of the relevant responses it produces. Fundamentally, it is vital to separate the model’s \emph{preferred outputs and actions} from its broader repertoire of \textit{knowledge and competencies}, enabling the construction of AI tools that are reliable, high-performing, and amenable to direction \citep{ouyang2022training}. Although current approaches often employ reinforcement learning (RL) to align language models with human preferences, we demonstrate that the RL-formulated objective commonly utilized by these approaches can in fact be achieved exactly by optimizing a simple binary cross-entropy objective—streamlining and reducing the complexity of the preference learning process.

In essence, prevailing techniques impart suitable behavioral patterns into LMs by drawing on thoughtfully constructed datasets reflecting human-preferred actions, thereby encouraging qualities such as safety and utility. This phase of learning from preferences is conducted subsequent to the initial extensive, unsupervised pre-training on a large-scale corpus of text. While directly supervising LMs with demonstrations of exemplary human responses constitutes the most straightforward strategy, the most effective methodologies to date leverage reinforcement learning from human (or AI) feedback (RLHF/RLAIF; \citep{christiano2017deep,bai2022constitutional}). RLHF methods first fit a reward model based on collected human preferences, then utilize RL to adapt the model’s output policy to maximize assigned reward, all while constraining it to avoid significant deviation from its initial parameters. Even though RLHF has enabled LMs to achieve impressive proficiency in dialogue and programming tasks, its pipeline is notably more intricate than conventional supervised fine-tuning. RLHF requires maintaining multiple LM instances and sampling from their policies during training, thereby increasing the computational demands substantially.

This work introduces a method for optimizing language models to align with human preferences without the need for explicit reward modeling or reinforcement learning. We present
\textit{{Direct Preference Optimization} (DPO)}, a novel approach that implicitly targets the same reward maximization objective under a KL-divergence constraint as standard RLHF methods, but with a design that is easy to implement and facilitates straightforward training. Conceptually, DPO updates function by enhancing the relative log likelihood of responses preferred by humans over those that are not, utilizing a dynamically determined, per-instance importance weighting. This weighting serves to avert the undesirable model collapse that can arise from a simple probability ratio approach. DPO, like prior work, utilizes a formal model of preference—such as the Bradley-Terry framework (\cite{bradley1952rankanalysis})—to quantify the correspondence between a candidate reward function and actual preference observations. Distinctly, whereas existing RLHF methodologies employ this preference model to train a reward model before optimizing a policy against it, DPO applies a variable substitution that enables the formulation of the preference loss directly in terms of the policy. Leveraging a collection of human preference annotations over model-generated responses, DPO can be trained with a straightforward binary cross entropy loss, thereby optimizing the policy to maximize an implicit reward function that best represents the preference data.

The principal contribution of this paper is Direct Preference Optimization (DPO), an uncomplicated algorithm for preference-based language model training that avoids RL. Experimental results demonstrate that DPO matches or surpasses the performance of established methods, including PPO-based RLHF, across a range of preference learning tasks such as sentiment control, summarization, and conversational modeling, utilizing language models scaling up to 6B parameters.

\section{Related Work}

As self-supervised language models grow in size, they demonstrate the capacity to address various tasks either in a zero-shot setting \citep{radford2019language} or using few-shot prompting strategies \citep{gpt3,megatron,chowdhery2022palm}. Nevertheless, their effectiveness on applied tasks and the degree to which their outputs align with user intentions can be greatly enhanced through a process of fine-tuning on curated instruction datasets, supplemented by human-generated completions \citep{mishra-etal-2022-cross,sanh2022multitask,chung2022scaling,thoppilan2022lamda}. This approach, known as `instruction-tuning,' facilitates broader generalization to novel instructions beyond those found in the training data and increases model practicality for end users \citep{chung2022scaling}. While instruction tuning has demonstrated substantial benefits, acquiring \textit{relative} human feedback—where annotators rank or compare responses—is typically more feasible than sourcing expert demonstration data. Consequently, subsequent studies have adapted large language models via fine-tuning on datasets that reflect human preferences, resulting in improved outcomes for translation \citep{kreutzer-etal-2018-reliability}, summarization \citep{stiennon2022learning,ziegler2020finetuning}, narrative generation \citep{ziegler2020finetuning}, and instruction-following abilities \citep{ouyang2022training,ramamurthy2023is}. The general approach consists of initially learning a neural reward function calibrated to these preference-labeled datasets using a preference aggregation scheme such as the Bradley-Terry model \citep{bradley1952rankanalysis}. Subsequently, the language model is refined to optimize this reward, typically by deploying reinforcement learning algorithms like REINFORCE \citep{williams1992reinforce}, proximal policy optimization (PPO; \cite{schulman2017proximal}), or related methodologies \citep{ramamurthy2023is}. Parallel work in this area utilizes instruction-tuned LLMs, guided by human feedback, to autonomously create supplementary synthetic preference data, focusing on properties like safety or benign behavior \citep{bai2022constitutional}; here, only lightweight human supervision is necessary, supplied in the form of a rubric used by the LLM for self-annotation. These methods unite two established research threads: one concerning reinforcement learning-based training of language models for diverse end goals~\citep{Ranzato2015SequenceLT,paulus2018a,wu2018learning}, and another dedicated to the study of human preference learning in general \citep{christiano2017deep,kupcsik2018learning}. Despite the promise shown by utilizing comparative human judgment, directly fine-tuning large language models via reinforcement learning is fraught with practical obstacles. The present work offers a theoretically-grounded alternative for optimizing models to respect relative preferences without reliance on RL.

Beyond the domain of language, the task of learning policies from preferences has garnered attention within both bandit and reinforcement learning frameworks, resulting in the development of several methodologies. In contextual bandit scenarios, where preferences or rankings over actions replace explicit reward signals, this is referred to as the contextual dueling bandit (CDB; \cite{yue2012karmed,dudik2015contextual}). Lacking absolute reward feedback, the theoretical foundations of CDBs adopt the concept of a \textit{von Neumann winner}, identifying a policy that is expected to outperform any competing policy in at least half of all matchups \citep{dudik2015contextual}. Distinctly, while CDBs receive preference feedback sequentially in an online manner, preference-based learning from human judgments typically operates with a predetermined dataset comprising offline preference-labeled action pairs \citep{yan2022human}. Likewise, \textit{preference-based RL} (PbRL) derives learning signals from binary preferences computed by an opaque `scoring' function instead of explicit rewards \citep{BusaFekete2014,ruiz2023dueling}. A variety of PbRL algorithms have been put forward—some capable of utilizing off-policy preference data—yet these strategies commonly follow a two-step procedure: first, they explicitly estimate the hidden scoring or reward function and subsequently optimize policies with respect to it \citep{jain2013learning,BusaFekete2014,christiano2017deep,sadigh2017active,kupcsik2018learning}. Contrasting these prior works, we introduce a unified, single-stage policy learning method that seeks to optimize the policy directly in accordance with the observed preferences.

\section{Preliminaries}\label{section:prelims}

We summarize the RLHF methodology as outlined in \citeauthor{ziegler2020finetuning} and expanded upon in subsequent works \citep{stiennon2022learning, bai2022training, ouyang2022training}. The pipeline consists of three main components: 1) supervised fine-tuning (SFT), 2) collection of preference data and reward modeling, and 3) optimization via reinforcement learning.

\textbf{SFT}: The RLHF process commences by applying supervised learning to a pre-trained language model, utilizing carefully curated datasets tailored to the downstream task (such as dialogue generation or summarization). This produces an SFT model, denoted $\pi^\text{SFT}$.

\textbf{Reward Modelling Phase}: In the second step, the fine-tuned SFT model is used to respond to prompts $x$, generating answer pairs $(y_1, y_2)\sim \pi^\text{SFT}(y \mid x)$. Human annotators then evaluate these responses, indicating their preference as $y_w\succ y_l \mid x$, with $y_w$ and $y_l$ representing the favored and less favored responses from the candidate pair, respectively. It is posited that these preferences are outcomes of some unobserved reward function $r^*(y, x)$. While several models exist to represent preferences, the Bradley-Terry (BT) model \cite{bradley1952rankanalysis} is frequently adopted (although the approach naturally accommodates broader ranking models such as Plackett-Luce \citep{plackett1975analysis, luce2012individual}, especially when rankings beyond pairs are available). Under the BT formulation, the human-generated preference probability distribution $p^*$ is expressed as follows:

Given a static dataset of pairwise comparisons $\mathcal{D}=\bigl\{x^{(i)}, y_w^{(i)}, y_l^{(i)}\bigr\}_{i=1}^N$ drawn from $p^*$, one can model a reward function $r_{\phi}(x, y)$ and estimate its parameters via maximum likelihood. Viewing this task as a binary classification problem, the corresponding negative log-likelihood loss is

\begin{equation}\label{eq:reward_model}
    \mathcal{L}_R(r_{\phi}, \mathcal{D}) = -\mathbb{E}_{(x, y_w, y_l)\sim \mathcal{D}}\bigl[\log \sigma(r_{\phi}(x, y_w)- r_{\phi}(x, y_l))\bigr]
\end{equation}

where $\sigma$ denotes the logistic sigmoid. In language modeling scenarios, $r_{\phi}(x, y)$ is typically seeded with the weights of the SFT model $\pi^\text{SFT}(y \mid x)$, supplemented by a linear layer on top of the terminal transformer layer to yield a single scalar score for the reward~\cite{ziegler2020finetuning}. To reduce the variance of the learned reward, previous works often normalize it such that $\mathbb{E}_{x,y\sim \mathcal{D}}\left[r_\phi(x, y)\right] = 0$ over all $x$.

\textbf{RL Fine-Tuning Phase}: In the subsequent RL fine-tuning phase, the policy receives signals from the learned reward. Building on prior work~\citep{jaques2017sequence, jaques2020human}, policy optimization is posed as

\begin{equation}\label{eq:RL}
\max_{\pi_{\theta}}  \mathbb{E}_{x\sim \mathcal{D}, y\sim \pi_{\theta}(y \mid x)}\bigl[r_{\phi}(x, y)\bigr] - \beta\mathbb{D}_{\textrm{KL}}\bigl[\pi_{\theta}(y\mid x)\mid \mid \pi_\text{ref}(y\mid x)\bigr],
\end{equation}

with $\beta$ controlling how much the policy $\pi_{\theta}$ is regularized towards the base reference policy $\pi_\text{ref}$ (typically initialized as $\pi^\text{SFT}$). The regularization term is crucial: it ensures that $\pi_\theta$ remains close to the data distribution for which the reward model is accurate, thereby maintaining sample diversity and preventing collapse onto narrow sets of responses. Because of the discrete outputs of language models, this objective lacks a straightforward gradient and thus is optimized with reinforcement learning. The standard method~\citep{ziegler2020finetuning, stiennon2022learning, bai2022training, ouyang2022training} is to shape the reward as ${r(x, y) = r_{\phi}(x, y) -\beta (\log \pi_{\theta}(y\mid x) - \log \pi_\text{ref}(y\mid x))}$, then maximize it using PPO~\cite{schulman2017proximal}.

\section{Direct Preference Optimization}\label{sec:DPO}

Recognizing the challenges inherent in scaling reinforcement learning to large language models, we aim to construct a streamlined policy optimization method that relies directly on observed preferences. In contrast to the prevailing RLHF pipeline—which separates reward model training from RL-based policy optimization—our approach takes advantage of a specific reward function parameterization. This parameterization admits a closed-form expression for its optimal policy, eliminating the need for an explicit RL fine-tuning loop. As described below, our principal insight is to utilize an exact analytical relationship between a reward function and its associated optimal policy, allowing us to reformulate loss functions from the reward space to the policy space. This change-of-variables formulation obviates the need to train an explicit, independent reward model, yet still adheres to preference modeling assumptions such as those encoded in the Bradley-Terry framework. Ultimately, the policy network jointly represents the language model as well as the underlying (implicit) reward.

\textbf{Derivation of the DPO objective.} The starting point is the reinforcement learning objective previously established in Eq.~\ref{eq:RL}, considering an arbitrary reward function $r$. Building upon existing literature~\citep{peters2007reinforcement, peng2019advantage, korbak2022reinforcement, go2023aligning}, it is evident that the solution optimizing the KL-regularized reward maximization problem in Eq.~\ref{eq:RL} is given by:

\begin{equation}\label{eq:op_policy}
    \pi_r(y\mid x) = \frac{1}{Z(x)}\pi_\text{ref}(y\mid x)\exp\left(\frac{1}{\beta}r(x, y)\right),
\end{equation}

where $Z(x) = \sum_{y} \pi_\text{ref}(y\mid x) \exp\left(\frac{1}{\beta} r(x, y)\right)$ represents the normalization factor, or partition function. For the comprehensive derivation, refer to Appendix \ref{app:derivation1}. In practice, even when $r$ is approximated by the MLE-estimated reward $r_{\phi}$ for the true reward $r^*$, direct computation of $Z(x)$ remains computationally challenging~\citep{korbak2022reinforcement, go2023aligning}, making this formulation impractical for large-scale usage. Nonetheless, it is possible to reformulate Eq.~\ref{eq:op_policy} by isolating the reward function as a function of the optimal policy $\pi_r$, the reference policy $\pi_\text{ref}$, and the unknown $Z(x)$. By applying the logarithm to both sides of Eq.~\ref{eq:op_policy} and rearranging, we get:

\begin{equation}\label{eq:main_eq}
    r(x,y) = \beta \log \frac{\pi_r(y\mid x)}{\pi_\text{ref}(y\mid x)} + \beta \log Z(x).
\end{equation}

This change of variables can be employed for the true reward $r^*$ and its optimal policy counterpart $\pi^*$. Importantly, the Bradley-Terry preference model relies solely on reward differences for two candidate completions: ${p^*(y_1 \succ y_2 \mid x) = \sigma(r^*(x, y_1) - r^*(x, y_2))}$. Inserting the reformulated expression for $r^*(x, y)$ from Eq.~\ref{eq:main_eq} into the preference model (Eq.~\ref{eq:bradley-terry}) removes the partition function term, enabling the human preference likelihood to be re-expressed exclusively in terms of $\pi^*$ and $\pi_\text{ref}$. Therefore, the RLHF-optimal policy $\pi^*$ under the Bradley-Terry structure obeys the relation:

\begin{equation}\label{eq:objective}
    p^*(y_1\succ y_2 \mid x) = \frac{1}{1 + \exp\left(\beta \log \frac{\pi^*(y_2\mid x)}{\pi_\text{ref}(y_2\mid x)} - \beta \log \frac{\pi^*(y_1\mid x)}{\pi_\text{ref}(y_1\mid x)}\right)}
\end{equation}

Full derivations are included in Appendix~\ref{app:derivation2}. While this outcome is derived using the Bradley-Terry model, comparable derivations can be made for broader families of preference models, such as the Plackett-Luce framework~\citep{plackett1975analysis, luce2012individual}; further details are in Appendix~\ref{app:plackett_luce_models}.

Having reformulated the human preference likelihood in terms of the optimal policy (rather than the reward function), we can directly set up a maximum likelihood objective for any parametric policy $\pi_\theta$. Mirroring the structure of the reward modeling objective (see Eq.~\ref{eq:reward_model}), this yields the policy learning objective:

\begin{equation}\label{eq:optimum_model}
    \mathcal{L}_\text{DPO}(\pi_{\theta}; \pi_\text{ref}) = -\mathbb{E}_{(x, y_w, y_l)\sim \mathcal{D}}\left[\log \sigma \left(\beta \log \frac{\pi_{\theta}(y_w\mid x)}{\pi_\text{ref}(y_w\

By adopting this alternative parameterization, we are effectively modeling an implicit reward, for which the optimal policy aligns with $\pi_\theta$. Furthermore, this approach can be interpreted as fitting a reparametrized Bradley-Terry model, endowing it with certain desirable theoretical guarantees, such as consistency given appropriate assumptions on the distribution of preference data \cite{bong2022generalized}. Additional theoretical analysis and comparison with related methodologies are provided in Section~\ref{sec:theory}.

\textbf{How does the DPO update operate?} To gain a mechanistic perspective on DPO, it is instructive to inspect the gradient of its loss function, $\mathcal{L}_\text{DPO}$. The gradient with respect to the model parameters $\theta$ is given by:

\begin{multline*}\label{eq:gradient}
    \nabla_\theta \mathcal{L}_\text{DPO}(\pi_\theta;\pi_\text{ref}) = \\ -\beta\mathbb{E}_{(x, y_w, y_l) \sim \mathcal{D}} \bigg[\underbrace{\sigma(\hat{r}_\theta(x, y_l) - \hat{r}_\theta (x, y_w))}_\text{assigns larger magnitude when the model ranks rewards incorrectly}\bigg[\underbrace{\nabla_\theta\log \pi(y_w \mid x)}_\text{encourages higher probability of $y_w$} - \underbrace{\nabla_\theta\log\pi(y_l \mid x)}_\text{reduces probability of $y_l$}\bigg]\bigg],
\end{multline*}

where $\hat{r}_\theta(x, y) = \beta \log \frac{\pi_\theta(y \mid x)}{\pi_\text{ref}(y \mid x)}$ defines the implicit reward based on the outputs of the language model $\pi_\theta$ and the reference model $\pi_\text{ref}$ (see also Section~\ref{sec:theory} for more details). Intuitively, this gradient expression operates to raise the likelihood of the preferred responses $y_w$ while lowering the likelihood of the less preferred $y_l$. Crucially, each training instance is weighted according to the degree to which the implicit reward model $\hat{r}_\theta$ mistakenly assigns a higher reward to the dispreferred completion, modulated by $\beta$. This reflects both the error in the ranking and the strength of the KL divergence constraint. Our empirical results highlight that such weighting is vital: omitting it—as in a simplistic variant of the method without this term—can lead the language model to undesirable behaviors or collapse (see Appendix Table~\ref{tab:unlikelihood_generations}).

\textbf{DPO overview.} The standard DPO workflow involves the following steps: 1) For each prompt $x$, generate candidate completions $y_1, y_2 \sim \pi_\text{ref}(\cdot \mid x)$ and use human feedback to produce a dataset of preference tuples $\mathcal{D} = \{x^{(i)}, y_w^{(i)}, y_l^{(i)}\}_{i=1}^N$. 2) The language model $\pi_\theta$ is then trained by minimizing $\mathcal{L}_\text{DPO}$, utilizing the chosen $\pi_\text{ref}$, dataset $\mathcal{D}$, and target $\beta$. In real-world applications, there is a preference to employ existing public datasets of human preferences, rather than collecting new completions and annotations. As such datasets are commonly generated using $\pi^\text{SFT}$, we select $\pi_\text{ref} = \pi^\text{SFT}$ whenever feasible. If $\pi^\text{SFT}$ cannot be accessed, we alternatively set $\pi_\text{ref}$ by maximizing the likelihood over the preferred completions $(x, y_w)$, specifically, $\pi_\text{ref} = \argmax_{\pi}\mathbb{E}_{x, y_w \sim \mathcal{D}}\left[\log \pi(y_w \mid x)\right]$. This approach reduces discrepancies between the ideal, but unobservable, reference distribution and the actual $\pi_\text{ref}$ employed within DPO. Additional technical details, including hyperparameter configurations, can be found in Appendix~\ref{app:implementation}.

\section{Theoretical Analysis of DPO} Here, we provide a deeper theoretical understanding of the DPO framework, clarify its benefits, and connect its properties to those of actor-critic algorithms commonly used in RLHF settings (e.g., PPO~\cite{schulman2017proximal}).

\label{sec:theory}

\subsection{Your Language Model Is Secretly a Reward Model} The DPO method unifies the learning of preferences and policy by leveraging a single maximum likelihood loss, rather than fitting a distinct reward function and performing policy optimization via RL. Notably, the DPO objective in Eq. \ref{eq:main_eq} is formally a Bradley-Terry model, parameterized by a reward function $r^*(x, y) = \beta \log\frac{\pi^*_\theta(y \mid x)}{\pi_\text{ref}(y \mid x)}$; optimizing the parameterized policy $\pi_{\theta}$ in this setting parallels optimizing a reward model as in Eq. \ref{eq:reward_model} with a reparametrization. This section formalizes the underlying theory of this equivalence, demonstrating that it does not limit the expressivity of the reward functions learned, and enables the exact recovery of the optimal policy. We initiate our analysis by formalizing an equivalence relation for reward functions.

\begin{definition}
Two reward functions $r(x, y)$ and $r'(x, y)$ are said to be equivalent if and only if ${r(x, y)-r'(x, y) = f(x)}$ for some function $f$.
\end{definition}

This relation clearly forms an equivalence relation, thereby dividing the set of all reward functions into equivalence classes. We present two supporting lemmas below:

\begin{lemma}\label{lemma:same_prefrence}
Within the framework of Plackett-Luce, and specifically the Bradley-Terry model, any pair of reward functions belonging to the same equivalence class will induce identical preference distributions.
\end{lemma}

\begin{lemma}\label{lemma:same_policy}
    Reward functions that are members of the same equivalence class result in the same optimal policy for the constrained RL setting.
\end{lemma}

The proofs of these lemmas are direct, and we provide them in Appendix \ref{app:lemma1}. The first lemma corresponds to a recognized under-specification in the Plackett-Luce model family \cite{plackett1975analysis}. This indeterminacy necessitates the adoption of extra identifiability conditions when attempting to derive guarantees for MLE estimations based on Eq. \ref{eq:reward_model} \cite{bong2022generalized}. The second lemma asserts that all reward functions within a given equivalence class determine the same optimal policy. Thus, our main focus is on identifying any reward function from the optimal equivalence class for the purposes of our objective. The following theorem, with proof detailed in Appendix~\ref{app:thm1}, formalizes this result:

\begin{theorem}\label{thm:main}
    Provided mild assumptions are satisfied, every reward class that is consistent with Plackett-Luce (or specifically, Bradley-Terry) models can be captured via the reparameterization ${r(x, y) = \beta \log \frac{\pi(y\mid x)}{\pi_\text{ref}(y\mid x)}}$ for some model $\pi(y\mid x)$ and a chosen reference model $\pi_\text{ref}(y \mid x)$.
\end{theorem}

\begin{sproof}
    Take any reward function $r(x, y)$, which gives rise to the corresponding optimal model $\pi_r(y \mid x)$ as specified by Eq. \ref{eq:op_policy}. We will demonstrate that there exists a representative of the equivalence class of $r$ that can be expressed in the given reparameterized form. Define the projection $f$ as
\begin{equation}
    f(r; \pi_\text{ref}, \beta)(x, y) = r(x, y) - \beta\log\sum_{y}\pi_\text{ref}(y\mid x)\exp\left(\frac{1}{\beta}r(x, y)\right)
\end{equation}
This operator $f$ acts by normalizing the reward through subtracting the logarithm of the partition function with respect to $\pi_r$. Because this normalization term depends solely on $x$, $f(r; \pi_\text{ref}, \beta)(x, y)$ still belongs to the equivalence class associated with $r(x, y)$. Substituting $r$ with the right side of Eq.~\ref{eq:main_eq}, which applies for any reward, yields $f(r; \pi_\text{ref}, \beta)(x, y) = \beta \log \frac{\pi_r(y\mid x)}{\pi_\text{ref}(y\mid x)}$. Therefore, the projection $f$ maps $r$ to a representative in the equivalence class of $r$ with the intended structure, confirming that the reparameterization introduces no loss of generality to the reward model.
\end{sproof}

An alternative interpretation of Theorem~\ref{thm:main} is that it characterizes the precise reward function selected by the DPO reparameterization within each equivalence class—namely, the function that satisfies

\begin{equation}\label{eq:lag_p}
     \sum_{y}\underbrace{\pi_\text{ref}(y\mid x)\exp\left(\frac{1}{\beta}r(x, y)\right)}_{=\pi(y\mid x)\text{, by Thm.~\ref{thm:main} reparam.}} = 1,
\end{equation}

which ensures that $\pi(y\mid x)$ is a proper probability distribution with non-negative values summing to one. Notably, Eq.~\ref{eq:lag_p} corresponds to the partition function associated with the optimal policy defined by the reward $r(x, y)$, as derived from Eq.~\ref{eq:op_policy}. The core insight underpinning the DPO algorithm is its introduction of constraints onto the otherwise unconstrained Plackett-Luce preference models (including the Bradley-Terry model), thereby maintaining the full expressive capacity of the reward class, while simultaneously rendering the optimal policy in Eq.~\ref{eq:op_policy} analytically manageable for every prompt $x$.

\subsection{Instability of Actor-Critic Algorithms} Our analytical framework also lends itself to the investigation of instability phenomena encountered with commonly employed actor-critic algorithms in RLHF, such as PPO. In alignment with the RLHF process, our focus is on the RL fine-tuning phase described in Section~\ref{section:prelims}. Connections may be drawn to the control as inference perspective \cite{levine2018reinforcement} as applied to the constrained reinforcement learning problem formalized in \ref{eq:RL}. Considering a parameterized policy $\pi_{\theta}(y\mid x)$, optimization seeks to minimize $\mathbb{D}_{\text{KL}}[\pi_{\theta}(y|x) \mid \mid \pi^*(y\mid x)]$, where $\pi^*$ denotes the optimal policy induced by $r_{\phi}(y, x)$ according to Eq.~\ref{eq:optimum_model}. Through algebraic manipulation, this yields the following optimization criterion:

\begin{equation}\label{eq:AC}
    \max_{\pi_{\theta}}\mathbb{E}_{\pi_{\theta}(y\mid x)}\bigg[\underbrace{r_{\phi}(x, y) -\beta\log\sum_{y}\pi_\text{ref}(y\mid x)\exp\left(\frac{1}{\beta}r_{\phi}(x, y)\right)}_{f(r_{\phi}, \pi_\text{ref}, \beta)} - \underbrace{\beta\log\frac{\pi_{\theta}(y\mid x)}{\pi_\text{ref}(y\mid x)}}_{\text{KL}}\bigg]
\end{equation}

The objective function under consideration aligns with those optimized in previous studies 
\citep{ziegler2020finetuning, stiennon2022learning, bai2022training, ouyang2022training}, utilizing the DPO-equivalent reward associated with the $r_{\phi}$ reward class. Within this framework, the normalization component of $f(r_{\phi}, \pi_\text{ref}, \beta)$ can be interpreted as the soft value function corresponding to the reference policy $\pi_\text{ref}$. Although this normalization does not modify the set of optimal solutions, its omission can result in elevated variance for the policy gradient estimator, leading to instability during training. One potential approach to handle this normalization is to incorporate a learned value function, yet such a method is often challenging to optimize effectively. Alternatively, earlier approaches have normalized the rewards by leveraging a human completion baseline, which effectively serves as a single-sample Monte Carlo approximation of the normalization term. Distinctly, the DPO reparameterization framework provides a reward function that obviates the need for any explicit baseline.

\section{Experiments}
Here, we present an empirical study to assess DPO’s effectiveness in learning policies directly from human preference data. We begin by conducting experiments in a controlled text-generation environment, examining DPO’s sample efficiency and its balance between reward maximization and KL-divergence minimization against established preference optimization algorithms, notably PPO. We subsequently extend our analysis to larger language models and more demanding RLHF scenarios, such as summarization and dialogue tasks. Our results indicate that, with minimal hyperparameter adjustment, DPO matches or surpasses robust baselines including RLHF with PPO, as well as outperforming strategies that select the highest-scoring trajectory among $N$ samples with respect to a trained reward model. Prior to detailing the outcomes, we outline the experimental methodology; comprehensive information is provided in Appendix~\ref{app:exp_details}.

\textbf{Tasks.} We conduct experiments across three open-ended text generation tasks. In each setting, algorithms derive a policy using a dataset of preference pairs, denoted as $\mathcal{D}=\bigl\{x^{(i)}, y_w^{(i)}, y_l^{(i)}\bigr\}_{i=1}^N$. The first task, \textbf{controlled sentiment generation}, involves $x$, which serves as the beginning portion of a movie review sourced from the IMDb dataset \cite{maas-EtAl:2011:ACL-HLT2011}; the task objective is to generate a continuation $y$ displaying positive sentiment. For a controlled comparison, we programmatically create preference pairs based on model outputs, using a pre-trained sentiment classifier that verifies $p(\text{positive}\mid x,y_w)>p(\text{positive}\mid x,y_l)$. The SFT baseline is established by further training GPT-2-large on positive reviews from the IMDb training split, up to convergence (refer to App~\ref{app:sentiment_details} for methodological details). In the \textbf{summarization} task, $x$ represents a Reddit forum post, while the goal is to synthesize a summary $y$ capturing the primary arguments of the post. Mirroring previous studies, we utilize the Reddit TL;DR summarization dataset \citep{volske-etal-2017-tl} along with the human preference dataset compiled by \citeauthor{stiennon2022learning}. Here, an SFT model—trained on Reddit summaries written by humans\footnote{\url{https://huggingface.co/CarperAI/openai_summarize_tldr_sft}}—serves as the base, and we adopt the TRLX RLHF framework \citep{leandro_von_werra_2023_7790115}. The human judgments were collected by \citeauthor{stiennon2022learning} on outputs produced by a comparable SFT model. The third task, \textbf{single-turn dialogue}, presents $x$ as a user's message, ranging from scientific inquiries to personal advice requests. The objective is to generate a relevant and helpful reply $y$; for this, we use the Anthropic Helpful and Harmless dialogue dataset \citep{bai2022training}, which contains 170,000 exchanges between people and an automated assistant, each concluding with two language model-generated replies and a human-chosen preference. Unlike other settings, a pre-trained SFT model is not accessible; therefore, we construct the SFT model by fine-tuning an existing language model exclusively on the responses favored by humans.

\textbf{Evaluation.} Our evaluation strategy incorporates two complementary methods. In the controlled sentiment generation setting, to rigorously assess how well each algorithm solves the constrained reward maximization problem, we chart each method’s trade-off curve between achieved reward and KL-divergence relative to the reference policy. Since we possess access to an oracle reward (i.e., a sentiment classifier), this frontier can be precisely determined. In contrast, real-world applications lack a known ground truth reward. Accordingly, we measure each algorithm’s \textit{win rate} compared to a baseline policy by leveraging GPT-4 to approximate human judgment for summary quality and response helpfulness in the summarization and single-turn dialogue tasks, respectively. For summarization, baseline performance is defined by the test set’s reference summaries; for dialogue, the baseline is given by the preferred human response in the evaluation data. While prior research has indicated that LMs may outperform traditional metrics as automated evaluators \citep{Chen2023ExploringTU}, we additionally validate the reliability of GPT-4 as a surrogate evaluator through a human study, as detailed in Sec.~\ref{sec:human-judgments}. Our results demonstrate that GPT-4’s assessments are strongly aligned with human judgments, with agreement rates equal to or surpassing inter-annotator consistency among humans.

\textbf{Methods.} Beyond DPO, we benchmark several established techniques for aligning language models with human preferences. For summarization, we implement zero-shot prompting with \textbf{GPT-J} \citep{gpt-j}; for dialogue, we apply 2-shot prompting using \textbf{Pythia-2.8B} \citep{biderman2023pythia}. Our evaluation also covers the \textbf{SFT} model and its variant \textbf{Preferred-FT}, which undergoes supervised fine-tuning on the preferred output $y_w$ generated either by the SFT model (in both sentiment and summarization settings) or by a generic LM (for the dialogue task). Another approach we assess is \textbf{Unlikelihood}~\citep{welleck2019neural}, a pseudo-supervised technique that trains the policy to assign higher probability to $y_w$ while explicitly \textit{reducing} the likelihood of $y_l$, with an optional scaling factor $\alpha\in[0,1]$ modulating the contribution of the unlikelihood loss. We further evaluate \textbf{PPO} \citep{schulman2017proximal}, where the reward function is trained from preference data, and include \textbf{PPO-GT}, which serves as an oracle utilizing the ground-truth reward in the controlled sentiment scenario. For sentiment, we employ two PPO-GT variants: a standard off-the-shelf implementation \cite{leandro_von_werra_2023_7790115} and an enhanced version incorporating reward normalization and refined hyperparameter selection (these improvements are also applied to PPO with learned rewards). Lastly, we incorporate the \textbf{Best of $N$} approach, wherein $N$ outputs are sampled from the SFT model (or from Preferred-FT in dialogue), and the candidate receiving the top score from the learned reward model is chosen. Although this method yields strong performance by decoupling reward modeling from PPO optimization, its computational demands render it infeasible for sizable $N$, as it necessitates generating $N$ outputs for every query during evaluation.

\subsection{How well can DPO optimize the RLHF objective?}

The standard KL-constrained reward maximization objective employed in most RLHF methodologies is designed to balance the pursuit of higher reward with the necessity to limit policy divergence from a fixed reference policy. Consequently, evaluating algorithm performance requires simultaneous consideration of both the attained reward and the extent of KL divergence; obtaining marginally greater rewards at the cost of significantly increased KL may not be advantageous. Figure~\ref{fig:frontier-tldr-main} presents the reward-KL tradeoff frontiers for different methods in the sentiment analysis context. For each algorithm, we perform several independent training sessions, varying a hyperparameter that modulates policy conservativeness in each session (for PPO, the target KL $\in\{3,6,9,12\}$; for DPO, $\beta \in \{0.05,0.1,1,5\}$; for unlikelihood, $\alpha\in\{0.05,0.1,0.5,1\}$; for preferred-FT, different random seeds are used). In total, 22 experiments are conducted across the methods. After every increment of 100 training steps, continuing until performance stabilizes, each learned policy is evaluated on a suite of held-out prompts. During evaluation, we measure the mean reward under the actual reward model and compute the average sequence-level KL\footnote{Specifically, this is the total KL-divergence summed over all time steps.} between the trained policy and the reference policy, i.e., $\text{KL}\left(\pi\mid \mid \pi_\text{ref}\right)$. The results reveal that DPO establishes the most favorable frontier, achieving the highest rewards while maintaining low KL values. This outcome is particularly striking for several reasons. Notably, both DPO and PPO attempt to optimize the identical objective, yet DPO does so with much higher efficiency; DPO's balance between reward and KL uniformly surpasses that of PPO. Moreover, DPO is able to realize a superior tradeoff curve even in comparison to PPO configured with access to ground-truth rewards (PPO-GT).

\subsection{Can DPO scale to real preference datasets?}
\label{sec:dpo-real-datasets}

We proceed by assessing DPO's capability for fine-tuning on real-world preference datasets, specifically focusing on summarization and single-turn dialogue. In the domain of summarization, standard automatic metrics like ROUGE frequently demonstrate weak alignment with human judgments~\citep{stiennon2022learning}. Previous research has established that language models fine-tuned with PPO on datasets reflecting human preferences tend to yield higher quality summaries. Our experimental protocol involves generating outputs on the TL;DR summarization dataset's test partition using various methods, then calculating the average win rate against ground-truth references. Completions for each approach are sampled at a range of temperatures between 0.0 and 1.0; results are presented in Figure~\ref{fig:frontier-tldr-main} (right). DPO, PPO, and Preferred-FT all operate by further fine-tuning the identical GPT-J SFT model\footnote{\url{https://huggingface.co/CarperAI/openai_summarize_tldr_sft}}. Our findings indicate that DPO attains a win rate near 61\% at zero temperature, surpassing PPO, which records roughly 57\% at its optimal sampling temperature. DPO also achieves a greater peak win rate than the best-of-$N$ baseline approach. Importantly, since we did not systematically adjust DPO’s $\beta$ hyperparameter, the observed performance might not represent the method's upper limit. In addition, DPO's results are far less sensitive to changes in temperature compared to PPO, which, at higher temperatures, exhibits performance that drops back to that of the underlying GPT-J model. The Preferred-FT method, by contrast, yields only marginal improvements over the SFT baseline. We further conduct direct human preference comparisons between DPO and PPO in Section~\ref{sec:human-judgments}, reporting that DPO outputs sampled at temperature 0.25 were favored in 58\% of cases over PPO samples generated at the same temperature.

For single-turn dialogue, we assess various approaches on a subset of the Anthropic HH dataset test split \citep{bai2022training}, focusing on examples with a single exchange between human and assistant. In evaluations using GPT-4, we calculate win rates for each approach by using the dataset’s preferred completions as ground truth references. Due to the absence of an established SFT model for this dataset, we initiate experiments from a pre-trained Pythia-2.8B model. A reference model is obtained via Preferred-FT, trained on selected completions to ensure generated outputs remain in-distribution, after which DPO training is conducted. Our analysis also incorporates comparisons with the best completion selected from 128 Preferred-FT samples (with performance saturating at this value, as demonstrated in Appendix Figure~\ref{fig:best-of-n}), as well as a 2-shot prompt configuration using the Pythia-2.8B base model. Across optimal temperatures for each method, DPO either matches or outperforms all baselines. Additionally, we consider an RLHF agent trained via PPO on the Anthropic HH dataset\footnote{\url{https://huggingface.co/reciprocate/ppo_hh_pythia-6B}} based on resources from a widely used implementation\footnote{\url{https://github.com/CarperAI/trlx/tree/main/examples/hh}}; however, we were unable to find prompt settings or sampling temperatures where it surpasses the performance of the base Pythia-2.8B model. Drawing from our findings in TL;DR and noting that both PPO and Best of 128 methods are optimized towards the same reward signal, we use Best of 128 as an approximate stand-in for PPO-level outcomes. In sum, DPO emerges as the sole method that achieves computational efficiency while consistently surpassing the performance of preferred completions in the Anthropic HH dataset, and delivers results comparable to or better than the much more resource-intensive Best of 128 method. Furthermore, Figure~\ref{fig:dialogue-main} illustrates that DPO attains peak performance in relatively few training steps.

\subsection{Generalization to a new input distribution}

To further investigate how PPO and DPO respond to shifts in input distribution, we apply the trained policies from our Reddit TL;DR summarization study to a new domain—specifically, news articles from the test set of the CNN/DailyMail corpus \citep{nallapati-etal-2016-abstractive}. For consistency, we use the top-performing sampling temperatures identified on TL;DR (0 and 0.25). Table~\ref{tab:ood} summarizes the outcomes. We measure the win rate using GPT-4 to compare model-generated summaries to the reference summaries, employing the same GPT-4 (C) prompt as in the Reddit TL;DR analysis, but modifying “forum post” to “news article.” On this out-of-distribution dataset, DPO consistently yields superior performance over PPO. These findings suggest that DPO achieves robust generalization comparable to, or better than, PPO—even without utilizing the unlabeled Reddit TL;DR prompts that augment PPO training.

\subsection{Validating GPT-4 judgments with human judgments}
\label{sec:human-judgments}
To assess the trustworthiness of GPT-4's summary evaluations, we conduct a human subject study drawing on outputs from the TL;DR summarization task and leveraging two GPT-4 prompting strategies. The \textbf{GPT-4 (S)} prompt is direct, asking only which summary more effectively communicates the main information from the post. In contrast, the \textbf{GPT-4 (C)} prompt adds an explicit requirement for conciseness, as we observed that under the simple prompt, GPT-4 tends to favor longer, potentially more redundant summaries relative to human preferences. Full prompt texts can be found in Appendix~\ref{app:prompts}. The evaluation includes three match-ups: the top (DPO, temp. 0.25), bottom (PPO, temp. 1.0), and intermediate (SFT, temp. 0.25) performing systems, each compared with PPO's greedy decoding (its best temperature). This sampling was designed to ensure coverage of a wide range of sample qualities. Results show that both GPT-4 prompts achieve agreement levels with humans comparable to human-human agreement (with multiple human judgments collected for DPO and PPO-1 comparisons only, due to limited availability of raters), supporting the use of GPT-4 as a reliable stand-in for human assessment. The \textbf{GPT-4 (C)} prompt generally yields win rates that more closely mirror human evaluations, motivating its use in reporting principal results in Section~\ref{sec:dpo-real-datasets}. Further methodological information, including interface details and the rater roster, is provided in Appendix~\ref{app:human-study}.

\section{Discussion}
Preference-based learning offers a robust and scalable approach to developing language models that are both effective and aligned with human values. In this work, we have presented DPO, a streamlined methodology for training language models from preference data, bypassing the need for reinforcement learning techniques. Unlike conventional approaches that frame preference learning within the standard RL paradigm to leverage established RL algorithms, DPO establishes a direct correspondence between language model policies and reward structures. This allows models to optimize for human preferences using a straightforward cross-entropy loss, eliminating the reliance on reinforcement learning frameworks and preserving general applicability. Remarkably, DPO achieves performance on par with, or superior to, prevailing RLHF methods, including those utilizing PPO, even with minimal hyperparameter adjustment. As such, DPO significantly simplifies the process of aligning language models through human feedback.

\textbf{Limitations \& Future Work.} Our findings prompt several avenues for continued research. One open question is how the DPO-trained policies fare in generalization to out-of-distribution data, especially in comparison to models trained with explicit reward signals. While preliminary experiments indicate comparable generalization to PPO-based alternatives, a more exhaustive evaluation is necessary. Further investigation is needed into whether employing self-labeling using DPO-trained models can similarly exploit unlabeled prompt datasets. Another subject for future exploration is the manifestation of reward over-optimization in the DPO context and whether the modest decline in performance observed in Figure~\ref{fig:dialogue-main}-right exemplifies this phenomenon. Moreover, our assessments are limited to models with up to 6B parameters; extending DPO to considerably larger, state-of-the-art architectures is a promising direction. In terms of evaluation methodology, we observe that GPT-4-derived win rates are sensitive to prompting, indicating that subsequent work could focus on optimizing the elicitation of reliable assessments from automated systems. Ultimately, DPO's utility is not confined to language modeling, as it holds potential for guiding generative models across various other domains.